% Image credits for the photographs and AI illustrations of Book 1 (Primary and Middle School), Hindi edition.
% Machine-readable license records live in images/book1/CREDITS.md; keep this
% file in sync with frontmatter/image-credits.tex (one \item per item there).
% Names of people, institutions, collections, works and licences are kept
% exactly as in the English credits (Latin script); only the frame and the
% descriptions are translated.
\chapter*{चित्रों का श्रेय}

इस पुस्तक के रेखाचित्र, मानचित्र और चार्ट मौलिक हैं; प्रकाशित आँकड़ों से बनाए गए मानचित्र और
चार्ट अपना स्रोत चित्र-परिचय में या इस सूची में बताते हैं। फ़ोटो उनके अपने-अपने मुक्त लाइसेंसों के
अंतर्गत, उनके लेखकों के प्रति आभार के साथ, इस्तेमाल किए गए हैं; हर फ़ोटो के पूरे स्रोत-लिंक और
लाइसेंस के पते पुस्तक के भंडार में \texttt{images/book1/CREDITS.md} में दिए गए हैं।

\bigskip
फ़ोटो और मौलिक रेखाचित्रों के साथ-साथ, कुछ चित्र AI से बनाए गए चित्रण हैं, जो पुस्तक के
संपादकों के लिखे निर्देशों से OpenAI के चित्र-निर्माण से बनाए गए, और शामिल करने से पहले उनकी
वैज्ञानिक शुद्धता जाँची गई। वे केवल सामान्य दृश्य दिखाते हैं, कभी कोई असली जगह, असली व्यक्ति
या असली नमूना नहीं। हर बनाए गए चित्र का पूरा निर्देश भंडार में
\texttt{images/book1/ai/PROMPTS.md} में दर्ज है।

\bigskip
\noindent फ़ोटो और आँकड़े, अध्याय के अनुसार:
\begin{itemize}\small\setlength{\emergencystretch}{2em}
  \item अध्याय~\ref{ch:g2:water-on-earth} (पृथ्वी पर पानी): Apollo~17 से देखी गई पृथ्वी, 7~दिसंबर 1972 (``The Blue Marble''); NASA, सार्वजनिक डोमेन (Wikimedia Commons के ज़रिए)। 100 वर्ग NOAA ETOPO1 उच्चावच से निकाले गए हैं; 100 बाल्टियाँ U.S. Geological Survey Water Science School की तालिका से (Shiklomanov, 1993 के अनुसार)।
  \item अध्याय~\ref{ch:g3:stones-pebbles-sand} (पत्थर, कंकड़ और रेत): ग्रेनाइट की सूखी पत्थर की दीवार; फ़ोटो: Tony Atkin, Geograph, CC BY-SA 2.0 (Wikimedia Commons के ज़रिए)। स्विट्ज़रलैंड में पत्थर की पटियों की छतें; फ़ोटो: Michael C. Rygel, CC BY-SA 3.0 (Wikimedia Commons के ज़रिए)।
  \item अध्याय~\ref{ch:g3:continents-and-oceans} (ग्लोब): दोनों ग्लोब और विश्व मानचित्र Natural Earth की तटरेखाओं से बनाए गए हैं (सार्वजनिक डोमेन)।
  \item अध्याय~\ref{ch:g4:clouds-rain-snow} (बादल, बारिश और हिम): Kansas के ऊपर कपासी मेघ (S.~Corfidi), पक्षाभ मेघों के टुकड़े (R.~F.~Kresge), पर्वत चोटियों के चारों ओर स्तरी मेघ की परत, Arizona के मरुस्थल के ऊपर एक तड़ित झंझा (B.~Hooke का संग्रह), एक बहुत बड़ा ओला; सभी NOAA Photo Library, सार्वजनिक डोमेन (Wikimedia Commons के ज़रिए)।
  \item अध्याय~\ref{ch:g4:volcanoes-earthquakes-first-look} (ज्वालामुखी और भूकंप: पहली झलक): Kilauea पर महासागर में गिरता लावा, USGS Hawaiian Volcano Observatory; Mount St.~Helens, 18~मई 1980, फ़ोटो: Austin Post, USGS; 1989 के Loma Prieta भूकंप से अपनी नींव से उखड़ा एक घर, फ़ोटो: J.~K.~Nakata, USGS; सभी सार्वजनिक डोमेन (Wikimedia Commons के ज़रिए)। मानचित्र Smithsonian Global Volcanism Program की होलोसीन ज्वालामुखियों की सूची और U.S. Geological Survey की भूकंप सूची का उपयोग करता है।
  \item अध्याय~\ref{ch:g4:world-physical-map} (विश्व का भौतिक मानचित्र): International Space Station से खींचे गए हिमालय और अमेज़न, NASA, सार्वजनिक डोमेन (Wikimedia Commons के ज़रिए)। भौतिक मानचित्र NOAA ETOPO1 उच्चावच से बनाया गया है, Natural Earth की नदियों, झीलों और तटरेखाओं के साथ (सार्वजनिक डोमेन)।
  \item अध्याय~\ref{ch:g5:water-shapes-the-land} (पानी ज़मीन को आकार देता है): International Space Station से खींची गई Mamoré नदी (ISS007-E-10797) और Terra उपग्रह के MODIS उपकरण से देखा गया नील नदी का डेल्टा, NASA, सार्वजनिक डोमेन (Wikimedia Commons के ज़रिए)।
  \item अध्याय~\ref{ch:g5:weather-and-climate} (मौसम और जलवायु): Sonoran मरुस्थल, U.S. National Park Service, Saguaro National Park, सार्वजनिक डोमेन (Wikimedia Commons के ज़रिए)। जलवायु चार्ट NOAA NCEI U.S. Climate Normals 1991--2020 से बनाए गए हैं (सार्वजनिक डोमेन)।
  \item अध्याय~\ref{ch:g5:natural-hazards} (प्राकृतिक संकट): International Space Station से Ed Lu द्वारा खींची गई Hurricane Isabel की आँख, NASA, सार्वजनिक डोमेन; Washington State में 2014 का भूस्खलन, Mark Reid, U.S. Geological Survey, सार्वजनिक डोमेन (दोनों Wikimedia Commons के ज़रिए)।
  \item अध्याय~\ref{ch:g6:earth-a-planet} (पृथ्वी, सौर मंडल का एक ग्रह): Mariner~10 से शुक्र (NASA/JPL-Caltech); Apollo~17 से पृथ्वी (NASA); Apollo~8 से \emph{Earthrise} (NASA, Bill Anders); Hubble Space Telescope से मंगल और बृहस्पति (NASA, ESA, the Hubble Heritage Team and A.~Simon); सभी सार्वजनिक डोमेन (Wikimedia Commons के ज़रिए)। 1990 में Voyager~1 से देखा गया \emph{Pale Blue Dot}, 2020 का संस्करण (NASA/JPL-Caltech, PIA23645, सार्वजनिक डोमेन)। ग्रहों के आँकड़े NASA planetary fact sheet से।
  \item अध्याय~\ref{ch:g6:latitude-longitude} (अक्षांश, देशांतर और समय-क्षेत्र): National Maritime Museum, Greenwich में Harrison की H4, फ़ोटो: Colonel Warden, CC BY-SA 3.0 (Wikimedia Commons के ज़रिए)। तटरेखाएँ और शहर Natural Earth से (सार्वजनिक डोमेन)।
  \item अध्याय~\ref{ch:g6:rocks-and-minerals} (चट्टानें और खनिज): ग्रेनाइट, बेसाल्ट, जीवाश्मयुक्त चूना पत्थर, बलुआ पत्थर, संगमरमर, शेल, क्वार्ट्ज़ का एक क्रिस्टल और मस्कोवाइट अभ्रक की एक परत, फ़ोटो: James St.~John, CC BY 2.0; स्लेट की एक छत, फ़ोटो: Øyvind Holmstad, CC0 (सभी Wikimedia Commons के ज़रिए)।
  \item अध्याय~\ref{ch:g6:inside-the-earth} (पृथ्वी के अंदर): 1932 में Inge Lehmann, फ़ोटो: Even Neuhaus, Royal Library of Denmark, सार्वजनिक डोमेन (Wikimedia Commons के ज़रिए)। परतें, घनत्व और तरंगों की चालें Preliminary Reference Earth Model से बनाई गई हैं (A.~M. Dziewonski and D.~L. Anderson, 1981)।
  \item अध्याय~\ref{ch:g6:water-underground} (ज़मीन के नीचे का पानी): Big Spring, Missouri, फ़ोटो: Johnston9494, CC BY-SA 4.0 (Wikimedia Commons के ज़रिए)। सरंध्रताएँ R.~C. Heath, \emph{Basic Ground-Water Hydrology}, U.S. Geological Survey (1983) के अनुसार; मीठे पानी के हिस्से USGS Water Science School से।
  \item अध्याय~\ref{ch:g7:plate-tectonics} (प्लेट विवर्तनिकी: चलती हुई पहेली): \emph{Cynognathus}, \emph{Lystrosaurus}, \emph{Mesosaurus} और \emph{Glossopteris} की जीवाश्म पट्टियों के साथ फिर से जोड़े गए दक्षिणी महाद्वीप; U.S. Geological Survey, \emph{This Dynamic Earth}, सार्वजनिक डोमेन।
  \item अध्याय~\ref{ch:g7:plate-tectonics}: Carrizo Plain को काटता हुआ San Andreas भ्रंश; फ़ोटो: Robert E. Wallace, U.S. Geological Survey, सार्वजनिक डोमेन।
  \item अध्याय~\ref{ch:g7:plate-tectonics}: प्लेट सीमाएँ P.~Bird (2003) के अनुसार, GeoJSON: H.~Ahlenius (ODC-By 1.0); तटरेखाएँ Natural Earth से (सार्वजनिक डोमेन); भूकंप U.S. Geological Survey की सूची से; ज्वालामुखी Smithsonian Global Volcanism Program से।
  \item अध्याय~\ref{ch:g7:volcanoes} (ज्वालामुखी): 17~मई और 10~सितंबर 1980 को Johnston Ridge से Mount St.~Helens, फ़ोटो: Harry Glicken, U.S. Geological Survey; Kilauea पर लावा का एक फव्वारा, 1983, और हवा से Mauna Loa, 1985, फ़ोटो: J.~D. Griggs, U.S. Geological Survey; Mount Pinatubo का उद्गार स्तंभ, 12~जून 1991, फ़ोटो: Dave Harlow, U.S. Geological Survey; सभी सार्वजनिक डोमेन (Wikimedia Commons के ज़रिए)। मानचित्र और गिनतियाँ Smithsonian Global Volcanism Program की होलोसीन ज्वालामुखियों की सूची का उपयोग करती हैं; प्लेट सीमाएँ P.~Bird (2003) के अनुसार।
  \item अध्याय~\ref{ch:g7:earthquakes} (भूकंप): 26~दिसंबर 2004 की सुनामी के बाद Sumatra के तट के पास एक गाँव, U.S. Navy फ़ोटो: Philip A. McDaniel; 1983 के Borah Peak भूकंप की भ्रंश कगारें, फ़ोटो: R.~E. Wallace, U.S. Geological Survey; Hebgen Lake में धँसी सड़क, 1959, फ़ोटो: I.~J. Witkind, U.S. Geological Survey; सभी सार्वजनिक डोमेन (Wikimedia Commons के ज़रिए)। गिनतियाँ और गहराइयाँ U.S. Geological Survey की भूकंप सूची का उपयोग करती हैं; तीव्रता की तालिका U.S. Geological Survey के Modified Mercalli पैमाने का संक्षिप्त रूप है।
  \item अध्याय~\ref{ch:g7:sediments-and-fossils} (अवसाद और जीवाश्म): Hermits Rest से Grand Canyon, NPS फ़ोटो: Kristen M. Caldon, सार्वजनिक डोमेन; Capitol Reef National Park में एक अमोनाइट, एक ट्राइलोबाइट और तरंग-चिह्न, फ़ोटो: James St.~John, CC BY 2.0; Dinosaur Ridge पर डायनासोर के पदचिह्न, फ़ोटो: Dexter Perkins, CC0 (सभी Wikimedia Commons के ज़रिए)।
  \item अध्याय~\ref{ch:g7:geological-time} (पृथ्वी के युग): James Hutton, Henry Raeburn का बनाया चित्र, सार्वजनिक डोमेन (Wikimedia Commons के ज़रिए)। समय-मापक्रम की आयु और रंग International Chronostratigraphic Chart (International Commission on Stratigraphy) के हैं।
  \item अध्याय~\ref{ch:g8:the-atmosphere} (वायुमंडल): Topeka, Kansas में मौसमी गुब्बारा छोड़ते हुए, फ़ोटो: Stephen Corfidi, NOAA/National Weather Service; International Space Station से सूर्यास्त के समय पृथ्वी का किनारा, NASA, ISS Expedition~23; दोनों सार्वजनिक डोमेन (Wikimedia Commons के ज़रिए)। रूपरेखाएँ \emph{U.S. Standard Atmosphere, 1976} (NOAA, NASA and U.S. Air Force) से निकाली गई हैं।
  \item अध्याय~\ref{ch:g8:winds-global-circulation} (पवनें और वैश्विक परिसंचरण): 26~जुलाई 2022 को GOES-16 उपग्रह से देखे गए अमेरिका महाद्वीप, NOAA, सार्वजनिक डोमेन (Wikimedia Commons के ज़रिए)। व्यापारिक पवनों के मानचित्र की तटरेखाएँ Natural Earth से (सार्वजनिक डोमेन)।
  \item अध्याय~\ref{ch:g8:fronts-and-weather-maps} (वायु राशियाँ, वाताग्र और मौसम मानचित्र): Terra उपग्रह से देखा गया Iceland के पास एक अवदाब, NASA/GSFC, MODIS Rapid Response Team, Jacques Descloitres, सार्वजनिक डोमेन (Wikimedia Commons के ज़रिए)। वायु राशियों के मानचित्र की तटरेखाएँ Natural Earth से (सार्वजनिक डोमेन)।
  \item अध्याय~\ref{ch:g8:the-oceans} (महासागर): गहराई नापने की तैयारी करता HMS \emph{Challenger}, 1872, Challenger अभियान की रिपोर्टों से एक उत्कीर्ण चित्र, W.~A. Herdman, \emph{Founders of Oceanography} (1923) में पुनर्मुद्रित; Aqua उपग्रह के MODIS उपकरण से देखी गई गल्फ़ स्ट्रीम, NASA Earth Observatory; दोनों सार्वजनिक डोमेन (Wikimedia Commons के ज़रिए)। समुद्री तल की परिच्छेदिकाएँ NOAA ETOPO1 उच्चावच से बनाई गई हैं; वलय और प्लव मानचित्रों की तटरेखाएँ Natural Earth से (सार्वजनिक डोमेन)।
  \item अध्याय~\ref{ch:g8:ice-on-earth} (पृथ्वी पर बर्फ़): 1990 और 2016 में Pedersen Glacier, फ़ोटो: M.~Tetreau और D.~Kurtz, U.S. National Park Service; Antarctic Peninsula के पास एक हिम शेल्फ़ का किनारा, NASA, Operation IceBridge (M.-J. Viñas); सभी सार्वजनिक डोमेन (Wikimedia Commons के ज़रिए)। समुद्री बर्फ़ का ग्राफ़ NSIDC Sea Ice Index, version~4 (F.~Fetterer and others, 2025, doi:10.7265/a98x-0f50) का उपयोग करता है।
  \item अध्याय~\ref{ch:g9:climate-past-and-present} (पृथ्वी की जलवायु, अतीत और वर्तमान): ग्रीनलैंड के GISP2 बर्फ़ कोर का एक हिस्सा, NOAA Paleoclimatology, सार्वजनिक डोमेन (Wikimedia Commons के ज़रिए)। अभिलेख NASA GISTEMP version~4, मौना लोआ का कार्बन डाइऑक्साइड अभिलेख (NOAA Global Monitoring Laboratory and Scripps Institution of Oceanography), EPICA Dome~C तापमान (J.~Jouzel and others, 2007), अंटार्कटिका के बर्फ़ कोरों का कार्बन डाइऑक्साइड संयोजन (B.~Bereiter and others, 2015) और NOAA Laboratory for Satellite Altimetry के समुद्र तल से बनाए गए हैं; आकलन के कथन IPCC Sixth Assessment Report, Working Group~I, Summary for Policymakers (2021) के हैं।
  \item अध्याय~\ref{ch:g9:natural-risks} (प्राकृतिक जोखिमों के साथ जीना): नेवाडो देल रुइज़ का संकट मानचित्र, T.~L. Wright and T.~C. Pierson, \emph{Living with Volcanoes}, U.S. Geological Survey Circular 1073 (1992) से; USCGC \emph{William Hart} के दल द्वारा जाँचा जाता एक सुनामी चेतावनी प्लव, U.S. Coast Guard, 2021; दोनों सार्वजनिक डोमेन (Wikimedia Commons के ज़रिए)।
  \item अध्याय~\ref{ch:g9:earth-resources} (पृथ्वी से मिलने वाले संसाधन): ASTER से देखी गई चुकीकामाटा खान, NASA/METI/AIST/Japan Space Systems and the U.S./Japan ASTER Science Team, 2007, सार्वजनिक डोमेन; Belmont County, Ohio में सड़क की कटान में कोयले की एक सीम, फ़ोटो: James St.~John, CC BY 2.0 (दोनों Wikimedia Commons के ज़रिए)। सांद्रण गुणक D.~Perkins and others, \emph{Mineralogy}, और S.~Earle, \emph{Environmental Geology} से हैं।
  \item अध्याय~\ref{ch:g9:landforms} (स्थलरूप): सेवन सिस्टर्स, फ़ोटो: Arild Vågen, 2017, CC BY-SA 4.0; दक्षिणी ब्राज़ील के तट पर अर्धचंद्राकार टिब्बे, International Space Station से NASA के एक अंतरिक्ष यात्री की फ़ोटो, 2017, सार्वजनिक डोमेन; Carlsbad Cavern में आश्चुताश्म, National Park Service फ़ोटो: Daniel Leifheit, 2018, सार्वजनिक डोमेन (सभी Wikimedia Commons के ज़रिए)।
  \item अध्याय~\ref{ch:g9:maps-of-the-earth} (पूरी पृथ्वी के मानचित्र): लेना नदी का डेल्टा, 2000 का Landsat चित्र, USGS EROS Data Center via NASA Visible Earth, सार्वजनिक डोमेन; एक GPS उपग्रह, NASA की कलाकार-कल्पना, सार्वजनिक डोमेन (दोनों Wikimedia Commons के ज़रिए)। विश्व मानचित्रों की तटरेखाएँ Natural Earth से (सार्वजनिक डोमेन)।
\end{itemize}
